% !TEX program = lualatex
% !TeX encoding = UTF-8
\PassOptionsToPackage{unicode}{hyperref}
\PassOptionsToPackage{bookmarksnumbered,bookmarksopen}{hyperref}
\documentclass[12pt,a4paper]{article}

% ============================================================
% 日本語設定 (LuaLaTeX + luatexja)
% ============================================================
\usepackage{luatexja}
\usepackage{luatexja-fontspec}

% ラテン文字フォント
\setmainfont{Times New Roman}[Ligatures=TeX]
\setsansfont{Arial}[Ligatures=TeX]
\setmonofont{Consolas}[Scale=MatchLowercase]

% 日本語フォント – Yu Gothic (Windows に標準搭載)
\setmainjfont{Yu Gothic}[UprightFont = * Regular, BoldFont = * Bold]
\setsansjfont{Yu Gothic}[UprightFont = * Regular, BoldFont = * Bold]
\setmonojfont{Yu Gothic}[UprightFont = * Regular, BoldFont = * Bold]

% ============================================================
% その他のパッケージ
% ============================================================
\usepackage{amsmath,amssymb,amsthm,mathtools}
\usepackage{geometry}
\usepackage{booktabs}
\usepackage{graphicx}
\usepackage{hyperref}
\usepackage{url}
\usepackage{xcolor}
\usepackage{tikz}
\usetikzlibrary{arrows.meta, positioning, calc}

\geometry{margin=1in}

% ============================================================
% YUCT マクロ
% ============================================================
\newcommand{\Keff}{K_{\mathrm{eff}}}
\newcommand{\kappac}{\kappa_c}
\newcommand{\eps}{\varepsilon}
\newcommand{\dint}{D+I\!\cdot\!R}
\newcommand{\YPSDC}{YPSDC}
\newcommand{\YUCT}{YUCT}

% ============================================================
% 定理環境（日本語）
% ============================================================
\newtheorem{theorem}{定理}
\newtheorem{definition}{定義}
\newtheorem{proposition}{命題}
\newtheorem{prediction}{予測}   % Prediction → 予測

% ============================================================
% ハイパーリンク設定
% ============================================================
\hypersetup{
	colorlinks=true,
	linkcolor=blue,
	citecolor=blue,
	urlcolor=blue,
}

\begin{document}
	
	\begin{titlepage}
		\begin{center}
			\vspace*{1cm}
			\Huge \textbf{付録 U: 調整効率に基づく重力通信 — \\ 因果律を破ることなく古典的限界を超える}
			
			\vspace{1cm}
			\Large Alexey V. Yakushev\textsuperscript{1}
			
			\vspace{0.5cm}
			\large \textsuperscript{1}Yakushev Research, \\
			YUCT Theory: \url{https://yuct.org/}\\
			YPSDC Protocol: \url{https://ypsdc.com/}\\
			
			\vspace{2cm}
			\textbf DOI: \url{https://doi.org/10.5281/zenodo.18444598}\\
			
			\vspace{1cm}
			
			\vspace{0cm}
			\large 
			本稿の短縮版は以前に発表されており、以下で入手可能です：\\ \url{https://doi.org/10.5281/zenodo.18362308}\\
			\vspace{1cm}
			2026年3月 \\ \large V.2.0.
			
			\newpage
			\vspace{2cm}
			\begin{abstract}
				\noindent
				我々は、ヤクシェフ統一調整理論（YUCT）に基づいた長距離通信の新しいアプローチを提案する。重力場を自然に存在するグローバルな辞書として扱うことにより、短いインデックス（重力摂動）が遠隔の観測者間で状態を調整できることを示す。その有効効率 \(\Keff\) は形式的に光速の限界を超えるが、因果律には違反しない。鍵は、オフラインでの辞書配布（事前に存在する重力場）とオンラインでのインデックス伝送（光速の重力波）の分離にある。重力波は物質を減衰なく透過し、電磁信号に対して大きな利点を提供する。我々は基礎方程式を導出し、達成可能な効率を見積もり、高精度重力計を用いた実験的検証を提案する。この研究は、遅延が距離ではなく局所処理によって決定される新しい種類の通信システムへの道を開く。
			\end{abstract}
			
			\vspace{0.5cm}
			\noindent \textbf{キーワード:} YUCT, 重力通信, 調整効率, YPSDC, 重力波, 因果律, 辞書, インデックス.
		\end{center}
	\end{titlepage}
	
	\tableofcontents
	\newpage
	
	\section{はじめに}
	\label{sec:intro}
	
	古典的通信システムは、電磁波（電波、レーザー）または光速 \(c\) 以下で伝播する他の信号に依存している。深宇宙ミッションでは、この根本的な制限が避けられない遅延を課す：火星への信号は4分から24分かかり、リアルタイム制御やデータ転送を著しく妨げる。量子テレポーテーションのような未来的な概念でさえ、\(c\) に制限された古典チャネルを必要とするため、この限界を回避できない。
	
	ヤクシェフ統一調整理論（YUCT）\cite{Yakushev2026} はパラダイムシフトをもたらす：それは \emph{データ伝送}（ビットの物理的転送）と \emph{状態の調整}（事前知識の整合）を区別する。中核となるアイデアは \textbf{二相プロトコル} \(\YPSDC\) である：オフライン相で共有辞書 \(D\) を配布し、オンライン相では短いインデックス \(\kappa\) のみを送信して \(D\) 内の対応するエントリを活性化する。調整効率
	\begin{equation}
		\Keff = \frac{H(\text{活性化されたアクション})}{H(\text{送信されたインデックス})}
	\end{equation}
	は、辞書が情報の大半を担うため、任意に大きくできる。重要なのは、インデックス自体は \(c\) より速く伝播しないこと；見かけ上の瞬時調整は事前に存在する辞書に由来する。
	
	本稿では、\emph{重力場}を自然で遍在する辞書として利用する可能性を探る。すべての質量を持つ物体は普遍的な重力相互作用に参加しており、それは連続的に分布した辞書と見なせる。質量を変調するなどして重力場の小さな摂動に情報を符号化することで、受信側が同じ重力辞書を用いてインデックスを復号できる。重力波は物質と非常に弱く相互作用するため、障害物を減衰なく透過する——これは電磁信号に対する決定的な利点である。さらに、重力場はすべての物体を先験的に結びつけるため、調整効率は原理的に光速の壁を因果律に違反せずに超えることができる。
	
	\section{YUCT の基礎}
	\label{sec:foundations}
	
	\subsection{\(\YPSDC\) プロトコル}
	\label{subsec:ypsdc}
	
	二人の観測者 \(A\) と \(B\) が、インデックス \(\kappa_i\) をアクション \(A_i\) に写像する辞書 \(\mathcal{D} = \{ (\kappa_i, A_i) \}\) を共有しているとする。辞書はオフラインで（おそらく亜光速で）配布される。オンライン通信中、\(A\) は \(\kappa_i\) を物理チャネルを通じて送信し、\(B\) はそれを受信して直ちに \(\mathcal{D}\) から \(A_i\) を取得する。調整の総時間は \(\kappa_i\) の伝送時間とルックアップ時間の和であり、後者は無視できるほど小さくできる。
	
	\subsection{調整効率}
	\label{subsec:keff}
	
	各アクション \(A_i\) が \(n\) ビットの情報を含み、インデックス \(\kappa_i\) が \(\ell\) ビットを持つとすると、\(\Keff = n/\ell\) である。一般に、
	\begin{equation}
		\Keff = \frac{H(\mathcal{A})}{H(\mathcal{K})},
	\end{equation}
	ここで \(H\) はシャノンエントロピーを表す。最適に圧縮された辞書では、\(\Keff\) は \(10^6\) 以上の値に達しうる（例えば現代の通信プロトコルにおいて）。
	
	\subsection{因果律と \(K_{\mathrm{eff}} > 1\)}
	\label{subsec:causality}
	
	\(\Keff > 1\) を光速超えの情報転送と解釈するのは単純すぎる。しかし、\(\Keff\) は活性化された情報と送信ビットの比率を測るものであり、ビット自体の速度ではない。インデックス \(\kappa\) は依然として \(v \le c\) で伝わる；余分な情報は事前に配布された辞書に由来する。辞書は以前に亜光速手段で確立されているため、因果律は保存される。したがって、\(\Keff\) は特殊相対性理論を破ることなく任意に大きくできる。
	
	\section{普遍的な辞書としての重力}
	\label{sec:gravity}
	
	\subsection{なぜ重力か？}
	
	重力場は、グローバルな辞書の理想的な候補となるいくつかのユニークな特性を持つ：
	\begin{itemize}
		\item \textbf{普遍性}：質量エネルギーを持つすべての物体が重力に参加する。場はどこにでも存在する。
		\item \textbf{透過性}：重力波は物質と非常に弱く相互作用するため、惑星、恒星、銀河を実質的に減衰や散乱なしに透過できる。
		\item \textbf{事前存在構造}：重力場はすでに質量分布を符号化しており、物体が存在する瞬間からすべての物体が共有する辞書と見なせる。
	\end{itemize}
	
	\subsection{辞書としての重力場}
	
	YUCT では、辞書は可能な状態の集合である。重力通信では、計量テンソル \(g_{\mu\nu}\)（またはその摂動）を辞書として解釈する。受信側は周囲の重力場（例えば地球の場、太陽の場など）を知っており、それからのずれを検出できる。送信側は局所的な重力場を変更する——例えば質量を動かしたり、形状を変えたり、共鳴を励起したりして、小さな摂動 \(\delta g_{\mu\nu}\) を作り出す。この摂動がインデックス \(\kappa\) として機能する。高感度重力計または干渉計を備えた受信側は \(\delta g_{\mu\nu}\) を測定し、共有辞書（既知の背景場）を用いて特定のアクションとして解釈する。
	
	\subsection{重力における \(\YPSDC\) の類推}
	
	\begin{itemize}
		\item \textbf{オフライン相}：背景重力場は宇宙内のすべての質量の分布によって確立される。この場は時間の始まりから自動的に「配布」される。
		\item \textbf{オンライン相}：送信側は局所的な摂動（インデックス）を作り出し、それは光速で（重力波として）外向きに伝播する。受信側は摂動を検出し、既知の背景と比較して符号化された情報を抽出する。
	\end{itemize}
	
	摂動は \(c\) で伝わるため、超光速信号伝播は存在しない。しかし、送信側と受信側の間の\emph{調整}、すなわち小さな摂動が潜在的に大量の情報を伝えるという事実は、非常に効率的（\(\Keff \gg 1\)）になりうる。
	
	\section{重力通信の数学的モデル}
	\label{sec:model}
	
	\subsection{信号符号化}
	
	送信側が事前に合意された符号に従って局所質量 \(M(t)\) を変調するとする。簡単のため、2値符号化を考える：\(\Delta M\) の変化が時間 \(\tau\) 継続すれば「1」、変化がなければ「0」を表す。距離 \(r\) での重力波振幅は近似的に
	\begin{equation}
		h(t) \sim \frac{G}{c^4 r} \frac{d^2}{dt^2} \left[ M(t) \right],
	\end{equation}
	ここで \(G\) はニュートン定数である。エネルギー束は極めて小さいが、現代の検出器（LIGO、将来の宇宙ベース観測所）は \(10^{-22}\) ものひずみを測定できる。
	
	\subsection{チャネル容量と調整効率}
	
	重力波のチャネル容量は検出器のノイズフロアと利用可能な帯域幅によって制限される。しかし、\emph{調整容量}は、各検出波形を可能なテンプレートのライブラリと照合できるため、より高い。辞書が \(N\) 個の可能なメッセージを含み、各メッセージが一意の波形テンプレートで記述されるとする。受信側は入力信号とすべてのテンプレートの相互相関をとる；インデックスは実質的にテンプレート識別子である。テンプレートが直交するように設計されていれば、受信イベントごとに伝えられるビット数は \(\log_2 N\) である。物理波を送るのに必要なエネルギーや時間は \(N\) に依存しない。したがって、
	\begin{equation}
		\Keff = \frac{\log_2 N}{\text{(物理信号の実効ビット長)}}.
	\end{equation}
	例えば、\(10^6\) 個のテンプレートがあり、各物理信号が1秒間のデータをサンプリングレート1 kHzで占める場合、生データサイズは \(1000\) ビットだが、伝えられる情報は \(\log_2 10^6 \approx 20\) ビットである。したがって \(\Keff \approx 20/1000 = 0.02\) となり、1より小さい。\(\Keff > 1\) を達成するには、テンプレートが運ぶ情報に比べて非常に短い必要がある。原理的には、特定の周波数に同調した共振重力アンテナを用いることができ、特定の周波数での短いバーストが特定の辞書エントリを活性化し、インデックスは実質的に周波数そのものとなる。周波数は高精度で測定できるため、識別可能な周波数の数は膨大になり、大きな \(\Keff\) をもたらす。
	
	\subsection{ノイズと誤差スケーリング}
	
	YUCT の普遍誤差法則 \cite{AppendixL} によれば、
	\begin{equation}
		\eps = \kappac \alpha \Keff^{-\beta}, \quad \beta = 0.67,
		\label{eq:universal-scaling}
	\end{equation}
	ここで \(\eps\) は誤認識の確率である。これにより、誤差が許容できなくなる前に \(\Keff\) がどれだけ大きくできるかに基本原理的限界が設定される。重力通信では、誤差率は検出器ノイズとテンプレートの不一致によって支配されるが、普遍法則は最適な \(\Keff\) が存在することを示唆する。
	
	\section{電磁通信との比較}
	\label{sec:comparison}
	
	\subsection{減衰と透過性}
	
	電磁波は物質によって遮断または散乱される：電波は電離層、水、または惑星内部によって吸収されうる；光信号は見通し線を必要とする。重力波は実質的に影響を受けずにすべてを通過する。これは、送信側と受信側が惑星、恒星、または他の障害物の反対側にいる場合でも重力リンクが動作できることを意味する。
	
	\subsection{電力要件}
	
	検出可能な重力波を生成するには巨大な質量と加速度が必要である。現在の技術では、天文距離での通信に十分な人工重力波を生成できない。しかし、太陽系内の惑星間距離では、将来の高エネルギー機械システム（例えば大型回転質量や核駆動発振器）によって必要なひずみが達成可能かもしれない。さらに、既存の自然源（パルサーなど）をビーコンとして利用すれば、それらはすでに一種の重力「キャリア」を提供しており、それを変調できる可能性がある。
	
	\subsection{遅延}
	
	電磁波と重力波はどちらも \(c\) で伝播するため、片道光時間は同一である。重力チャネルの利点は速度ではなく \textbf{効率} にある：辞書がすでに存在するため、短いインデックスで複雑なメッセージを伝えることができ、事実上情報を圧縮する。これは遅延を減らさないが、必要な帯域幅とビットあたりのエネルギーを削減する。
	
	\subsection{調整効率の可能性}
	
	無線通信では、最大 \(\Keff\) はチャネル容量と変調の複雑さによって制限される。重力波では、辞書（背景重力場）が非常に豊かで安定しているため、潜在的な \(\Keff\) ははるかに高くなる可能性がある。例えば、地球–月系の重力ポテンシャルは、数百万の異なる状態（異なる軌道構成）を持つ辞書として機能しうる。小型衛星の軌道を変調することで、遠方の観測者が暦を知っていれば高効率で復号できる情報を符号化できる。
	
	\subsection{重力チャネルの非対称性}
	\label{subsec:asymmetry}
	
	電磁通信と重力通信の根本的な違いは、送信機と受信機の要件の非対称性にある。レーザーリンクが正確な指向と入射ビームに合わせた光検出器を必要とするのに対し、レーザー干渉計などの重力波検出器は全方向に動作する。それは常に時空ひずみを監視し、特徴的な信号（「コード」）を待つ——光源の方向を知る必要はない \cite{LIGO}。これにより、重力受信は指向と追跡の点で本質的に単純になる。
	
	しかし、検出可能な重力波を生成するには膨大なエネルギーと質量変位が必要である。地上局は原理的に大規模な共振質量や軌道変調器（例えば km スケールの干渉計 LIGO、または将来の宇宙ベース観測所）をホストできる。一方、小型宇宙船は惑星間距離で検出可能なひずみを生成するのに必要な質量と電力を欠いている。したがって、リンクは本質的に \textbf{一方向} である：強力な地上（または大型軌道）送信機から、単に聴くだけの受動的受信機へ。双方向通信は、宇宙船に同様に大規模な送信機を搭載するか（現状では非現実的）、または核源を動力とする高周波共振空洞のような根本的に新しいタイプの重力波エミッターを必要とするが、それはまだ推測の域を出ない。
	
	この非対称性は、多くの応用においてチャネルの価値を損なわない：深宇宙探査機は強力な搭載送信機を必要とせずに複雑なコマンド（インデックス）を受信でき、テレメトリは従来の無線で送り返せる。したがって、重力ダウンリンクは高効率ブロードキャストチャネルとして機能し、既存のアップリンクを補完する。
	
	将来の高エネルギー機械システムの発展や、自然発生の重力波源（例えば変調されたパルサー）をビーコンとして利用することで、いつかは対称リンクが可能になるかもしれない。当面は、一方向重力通信の実現可能性を実証し、制御された実験でその調整効率 \(K_{\mathrm{eff}}\) を測定することに重点を置くべきである。
	
	\section{信号伝播と重力ダイナミクスにおける調整場密度の役割}
	\label{sec:coord-density}
	
	\subsection{時空幾何の修飾子としての調整場密度}
	
	YUCT フレームワークでは、重力は単なる空虚時空の幾何学的特性ではなく、基礎となる調整場 \(\Psi_{MN}\) の顕現である。この場の密度 \(\rho_K = \langle \Psi_{MN}\Psi^{MN} \rangle\) は、情報の流れに対する屈折率に類似した役割を果たす。調整密度が高い領域（例えば大質量天体の近傍、コヒーレント量子系、または初期宇宙の相転移中）では、インデックスの伝播を支配する有効計量が修正される。
	
	完全な19次元ラグランジアン（付録 A）とその結合項（例えば共鳴的インターセクター結合を記述する \(\mathcal{L}_{329}\)）から出発して、調整場が創発計量に寄与する有効な4次元作用を導出できる：
	\begin{equation}
		g_{\mu\nu}^{\mathrm{eff}} = g_{\mu\nu} + \alpha \, \rho_K \, T_{\mu\nu}^{\mathrm{coord}},
	\end{equation}
	ここで \(\alpha\) は無次元定数、\(T_{\mu\nu}^{\mathrm{coord}}\) は調整場のエネルギー・運動量テンソルである。結果として、インデックス（場の小さな摂動）の速度はもはや厳密に \(c\) ではなく、
	\begin{equation}
		v_{\mathrm{index}} = c \left(1 + \beta \, \rho_K + \cdots \right),
	\end{equation}
	となる。ここで \(\beta\) は重力セクターと情報セクターの間の結合強度 \(\kappa_{0,103}\) に関連する別の定数である。この効果は通常の条件下では極めて小さいが、極端な天体物理学的または宇宙論的設定では測定可能になる可能性がある。
	
	\subsection{共鳴増強と好ましいチャネルの存在}
	
	\(\mathcal{L}_{329}\) の非線形項（例えば \(-\epsilon \Phi^3\) 項）の存在は、特定の周波数 \(\omega_{mn}\) での信号のパラメトリック増幅を可能にする。\(\rho_K\) が大きい領域では、これらの共鳴が「調整の川」——すなわちインデックスが最小の損失で伝播し、調整場自体からエネルギーを引き出すことによる増幅さえも伴う好ましい方向または周波数——を創り出すことができる。これはレーザー利得媒体に類似しているが、ここでは増幅が時空の幾何学的構造内で発生する。
	
	重力通信リンクにとって、これは効率 \(\Keff\) が辞書サイズだけでなく局所的な調整密度によっても決定されることを意味する。送信側が経路に沿った調整場の固有周波数に一致する共鳴周波数 \(\omega_{\mathrm{res}}\) で質量を変調できれば、信号は優先的にチャネリングされ、光源での必要な電力が大幅に低減される。同じ周波数に同調された受信側は共振検出器として機能し、背景ノイズからインデックスを抽出する。
	
	\subsection{信号伝播に関する予測}
	
	上記の考察から、既に Sect.~\ref{subsec:asymmetry} に挙げたものを補完する三つの検証可能な予測を得る：
	
	\begin{prediction}[重力波速度の方向依存性]
		調整場密度 \(\rho_K\) が等方的でない場合（例えば大規模構造や優先的なフレームのため）、重力波（したがって任意の重力インデックス）の速度は小さな方向依存性を示すはずである。異なる空方向からの重力波信号の到着時刻を相関させることで、\(\Delta v/c \sim 10^{-18}\)–\(10^{-20}\) のオーダーの異方性が明らかになる可能性があり、これは将来の第三世代検出器ネットワーク（アインシュタイン望遠鏡、コズミックエクスプローラー）の到達範囲内である。
	\end{prediction}
	
	\begin{prediction}[ミリヘルツ帯の共鳴周波数]
		\(\mathcal{L}_{329}\) の非線形項 \(\epsilon \Phi^3\) は、ミリヘルツ帯（銀河構造の調整場が固有モードを持つ可能性がある）に狭い共鳴周波数が存在することを予測する。文献~\cite{Barontini2025} で提案されているような超安定光学キャビティを用いて、二つの離れた検出器の相互相関でこれらの共鳴を探索すれば、既知の天体物理学的源では説明できない周波数 \(\omega_{mn}\) にピークが現れるはずである。
	\end{prediction}
	
	\begin{prediction}[調整場による人工信号の増幅]
		人工重力波源（例えば大型回転質量）が局所調整場（例えば地球–月系）の共鳴周波数に同調されている場合、遠方検出器での受信信号強度は標準の四重極公式による予測を因子 \(\sim Q\) だけ超えるはずである。ここで \(Q\) は共鳴の品質因子（潜在的に \(10^3\)–\(10^6\)）である。これは調整場増幅の直接的な実証となるであろう。
	\end{prediction}
	
	\subsection{普遍スケーリング則との関連}
	
	普遍誤差スケーリング則 \(\eps = \kappac\alpha\Keff^{-\beta}\)（式~\ref{eq:universal-scaling}）はインデックスの伝播にも適用される。\(\rho_K\) が増強された領域では、与えられたインデックス長に対する実効的な \(\Keff\) が増加し、誤差確率が減少する。これにより、因果律を破ることなく非常に高い調整効率を達成するための自然なメカニズムが提供される：辞書はすでに存在し、インデックスはその配信を積極的に支援する媒質を伝わる。
	
	\section{提案する実験的検証}
	\label{sec:experiment}
	
	\subsection{セットアップ}
	
	地球上の離れた場所、例えば1000 km離れた二箇所に高感度重力計（例えば超伝導重力計または原子干渉計）を設置する。一方の場所で、大型質量（例えば数トンの振り子）を所定のパターンで変調する。重力信号は \(c\) で伝播し、他方の場所で検出される。検出された信号を既知のパターンと相関させることで、実効的な調整効率を測定し、理論的予測と比較できる。
	
	\subsection{予測される結果}
	
	比較的単純な辞書（例えば振幅コードのセット）であっても、\(\Keff\) は同じエネルギー消費に対する古典的限界を超えると予想される。実験は普遍誤差法則も検証する：コード数を増やす（\(\Keff\) を増やす）と、誤差率は \(\eps \propto \Keff^{-0.67}\) に従うはずである。
	
	\subsection{課題}
	
	主な課題は人工重力波の振幅が極めて小さいことである。1000 km の距離で、1000トンの質量が1 Hz で1 m の振幅で振動すると、ひずみは約 \(10^{-23}\) のオーダーになり、現在の検出限界の限界にある。将来の重力波検出器（例えばアインシュタイン望遠鏡）は必要な感度を達成するかもしれない。
	
	\section{考察}
	\label{sec:discussion}
	
	YUCT に基づく重力通信の概念は、従来の思考から根本的な逸脱を提供する。事前に存在する重力場を辞書として活用することで、因果律を厳格に守りながら光速を超えるように見える調整効率を達成できる。実用的実装は計り知れない技術的ハードルに直面するが、理論的枠組みは健全で検証可能である。
	
	成功すれば、そのようなシステムは太陽系全体にわたる瞬時の（ユーザーから見た）通信を可能にするかもしれない。なぜなら、遅延は距離ではなくインデックスの生成と検出の時間によって支配されるからである。さらに、重力波はすべての物質を透過するため、地下深くや海洋下を含むどこからでも通信が可能になる。
	
	\section{結論}
	\label{sec:conclusion}
	
	我々は、ヤクシェフ統一調整理論に基づく重力通信の理論的枠組みを提示した。重力場を自然に存在するグローバルな辞書として扱うことで、短いインデックス（重力摂動）が遠隔の観測者間でアクションを調整できることを示した。その効率 \(\Keff\) は光速によって制限されない。インデックス自体は依然として \(c\) で伝播し、因果律を保存する。重力波は透過性と普遍性というユニークな利点を提供し、将来の深宇宙通信システムにとって魅力的である。我々は既存技術を用いた潜在的な実験的検証を概説した。この研究は、重力物理学と情報理論の両方に新たな方向性を開き、星間通信の未来に深い影響を与える。
	
	\section*{謝辞}
	
	著者は、YUCT セミナーの参加者に有益な議論に感謝する。
	
	\begin{thebibliography}{99}
		
		\bibitem{Yakushev2026}
		A. V. Yakushev, \emph{Yakushev Unified Coordination Theory (YUCT)}, Zenodo, \url{https://doi.org/10.5281/zenodo.18444599} (2026).
		
		\bibitem{AppendixL}
		A. V. Yakushev, \emph{Appendix L: Fractal Coordination Error Scaling — A Universal Law from DNA to Cosmology}, Zenodo (2026).
		
		\bibitem{AppendixA}
		A. V. Yakushev, \emph{Appendix A: Practical Guide to Using YUCT V35.0 for Interdisciplinary Modeling and Predictions}, Zenodo (2026).
		
		\bibitem{LLR-IfE}
		J. M. et al., "Lunar Laser Ranging: A Continuing Legacy of the Apollo Program", \emph{Science} \textbf{265}, 482–490 (1994).
		
		\bibitem{LIGO}
		B. P. Abbott et al. (LIGO Scientific Collaboration), "Observation of Gravitational Waves from a Binary Black Hole Merger", \emph{Phys. Rev. Lett.} \textbf{116}, 061102 (2016).
		
		\bibitem{Barontini2025}
		F. Barontini et al., "Detecting millihertz gravitational waves with optical cavities", \emph{Phys. Rev. Lett.} (to be published) (2025).
		
	\end{thebibliography}
	
\end{document}